\documentclass[11pt,letterpaper]{amsart}
\IfFileExists{lmodern.sty}{\usepackage[T1]{fontenc}\usepackage{lmodern}}{}
\usepackage{mathtools,amssymb,booktabs}
\usepackage[margin=1.32in,includehead,headsep=17pt]{geometry}
\usepackage[hidelinks]{hyperref}
\hypersetup{
  pdftitle={Elliptic Quotient Spectra under Complex Multiplication},
  pdfauthor={Crest--Root Calculus},
  pdfsubject={Polarized isogenies, finite kernels, and optimal elliptic quotients}
}
\newtheorem{theorem}{Theorem}[section]
\newtheorem{corollary}[theorem]{Corollary}
\newtheorem{lemma}[theorem]{Lemma}
\newtheorem{proposition}[theorem]{Proposition}
\theoremstyle{definition}
\newtheorem{definition}[theorem]{Definition}
\numberwithin{equation}{section}
\newcommand{\Q}{\mathbb Q}
\newcommand{\Z}{\mathbb Z}
\newcommand{\C}{\mathbb C}
\newcommand{\OO}{\mathcal O}
\newcommand{\NN}{\mathrm N}
\newcommand{\av}{\operatorname{Avg}}
\newcommand{\vol}{\operatorname{vol}}
\newcommand{\abs}[1]{\lvert #1\rvert}

\title[Elliptic quotient spectra under complex multiplication]{Elliptic Quotient Spectra under Complex Multiplication}
\author{Crest--Root Calculus}
\date{26 November 2026}
\subjclass[2020]{Primary 14H30, 14K22. Secondary 11G50, 11H06, 14K12.}
\keywords{elliptic subcovers, complex multiplication, polarized isogenies,
finite kernels, Hermitian forms, height counting}

\begin{document}
\begin{abstract}
Let a polarized complex abelian variety be isogenous to a power of an
elliptic curve with complex multiplication. We express the degree of
each elliptic subvariety as a Hermitian norm divided by the order of an
intersection with the finite isogeny kernel. The relevant norm is
normalized by the ideal class of the direction. Counting primitive
vectors in ideal lattices gives an asymptotic with an explicit leading
constant. The finite correction is a product of local averages of
powers of the kernel intersection orders. For two genus-two Jacobians
with the same rational polarized isogeny data over a common number
field, the leading constants are in the ratio \(3:4\). For the
non-product \(j=1728\) two-torsion gluing we give the multiplicity of
every optimal degree as a divisor formula and obtain a remainder of
order \(X\log(2+X)\). We give two pairs of explicit genus-two curves
over quartic fields whose optimal quotient counts agree in every degree
prime to three. The first differences occur at degrees three and six,
respectively. Their identifications with the corresponding Humbert
lattices are certified from the equations.
\end{abstract}
\maketitle

\section{Introduction}

For a smooth curve \(C\) of genus two over \(\C\), an elliptic quotient
is a nonconstant morphism \(f:C\to E\) to an elliptic curve. Call it
\emph{optimal} if it does not factor through a nontrivial isogeny onto
\(E\). We identify maps that differ by a translation or an isomorphism
of their targets. An optimal class corresponds to an elliptic
subvariety of the Jacobian \(J(C)\). Its degree is the degree of the
canonical polarization on that subvariety
\cite{KaniI,KaniII,Guerra19}. We write \(m_C(d)\) for the number of
classes of degree \(d\), and \(N_C(X)=\sum_{d\le X}m_C(d)\).

These counts record integral data. An isogeny identifies elliptic
directions on its source and target, but can change the degree of a
direction by a factor that depends on how its torsion meets the
isogeny kernel. Two kernels of equal order can therefore produce
different degree distributions. The question treated here is how to
compute this change for every direction, including when the relevant
complex multiplication field has nontrivial class group.

Our main result, Theorem~\ref{thm:height}, gives a counting law for
every polarized variety isogenous to a CM elliptic power in
characteristic zero. The formula uses the determinant of the
pulled-back Hermitian form, the ideal class number and a finite
average on torsion. The theorem applies to a specified isogeny. Any
CM elliptic power has such a presentation using an elliptic curve
with maximal order, since elliptic curves in a fixed CM isogeny
class are isogenous to \(\C/\OO\) for the full ring of integers.

The finite factor has geometric consequences. In
Section~\ref{sec:minus7}, two gluings of a single elliptic square have
the same pulled-back polarization and the same rational symplectic
Tate modules over a common field. They have different exact
multiplicities and different leading constants. Section~\ref{sec:gaussian}
computes every multiplicity in the Gaussian case by reducing to the
four-squares formula, followed by ideal M\"obius inversion. This gives
an independent description of the CM lattices in the accompanying
certificates.

Kani's refined Humbert invariant already identifies covers using
primitive representations of an intrinsic quadratic form
\cite{KaniI,KaniII}. Guerra treats the degree of elliptic subvarieties,
counts in products and comparison under polarized isogenies
\cite[Sections 6 and 7]{Guerra19}. Ideal classes, fundamental
domains for units and M\"obius inversion also occur in Schanuel's
projective height count \cite{Schanuel}. Our calculation evaluates the
restricted isogeny degree as a function of the primitive residue
class and inserts that function into the leading constant. The
examples use this term to distinguish spectra in a fixed polarized
rational isogeny class.

\section{Directions and the counting theorem}

Fix an imaginary quadratic field \(F\) with ring of integers \(\OO\),
discriminant \(D_F<0\), class number \(h_F\), and \(w_F=\abs{\OO^\times}\).
Choose the complex embedding and set \(\Delta=\sqrt{\abs{D_F}}/2\).
Thus \(\Delta\) is the covolume of \(\OO\) in \(\C\).
Let \(E=\C/\OO\) and
\[
  \Psi:E^r\longrightarrow A
\]
be an isogeny onto a polarized abelian variety \((A,\mathcal L)\) of
dimension \(r\ge2\). Let \(K=\ker\Psi\) and choose a positive integer
\(n\) killing \(K\). With the product principal polarization on \(E^r\)
as reference, \(\Psi^*\mathcal L\) is represented by a positive
definite Hermitian matrix \(H\) with entries in \(\OO\). Its norm is
\(Q(v)=v^*Hv\).

For a nonzero integral ideal \(\mathfrak a\), call \(v=(v_1,\dots,v_r)\)
primitive in \(\mathfrak a^r\) if \(\sum_i v_i\OO=\mathfrak a\).
Choose one ideal for each ideal class. Define
\[
 E_{\mathfrak a}=\C/\mathfrak a^{-1},\qquad
 \iota_v:E_{\mathfrak a}\longrightarrow E^r,\quad
 z\longmapsto(v_1z,\ldots,v_rz).
\]
The map is well-defined and its kernel is zero, since
\(\{z\in\C:\ v_i z\in\OO\text{ for all }i\}=\mathfrak a^{-1}\).
The integers \(n\) and
\begin{equation}\label{eq:ka}
 k_{\mathfrak a}(\bar v)=
 \#\{P\in E_{\mathfrak a}[n]:\ \iota_v(P)\in K\},
 \qquad \bar v=v\bmod n\mathfrak a ,
\end{equation}
give a finite calculation. The right side of \eqref{eq:ka} only uses
\(\bar v\), since multiplying the difference of two lifts by
\(P\in E_{\mathfrak a}[n]\) gives zero coordinatewise in \(E\).
For the unit ideal write \(k=k_{\OO}\). Call a residue
\(\bar v\in(\OO/n\OO)^r\) admissible if, for each prime ideal
\(\mathfrak p\mid n\OO\), its coordinates are not all zero modulo
\(\mathfrak p\).

\begin{theorem}\label{thm:height}
Every elliptic subvariety of \(A\) is represented by one primitive
\(v\in\mathfrak a^r\) for one of the chosen \(\mathfrak a\), up to
multiplication of \(v\) by a unit. Its degree is
\begin{equation}\label{eq:degree}
 d(v)=\frac{Q(v)}{\NN(\mathfrak a)\,k_{\mathfrak a}(\bar v)}.
\end{equation}
Let \(N_A(X)\) count elliptic subvarieties of polarization degree at
most \(X\). Then
\begin{align}
 N_A(X)&=c_A X^r+O_{F,H,n}(X^{r-1/2}), \label{eq:main}\\
 c_A&=\frac{h_F}{w_F\,\Delta^r\,\zeta_F(r)}
        \frac{\pi^r}{r!\det H}\,
        \av_{\bar v\ {\rm admissible}} k(\bar v)^r. \label{eq:constant}
\end{align}
The last average factors over the rational primes dividing \(n\).
\end{theorem}

\begin{proof}
An isogeny takes connected elliptic subvarieties bijectively to
connected elliptic subvarieties, with inverse given by the identity
component of the inverse image. In \(E^r\) these subvarieties
correspond to \(F\)-lines in \(F^r\). Indeed, an elliptic subvariety
projects nontrivially to some copy of \(E\), hence is isogenous to
\(E\). After choosing an isogeny from \(E\), its coordinates are
endomorphisms of \(E\), and the resulting line is independent of the
chosen isogeny. Scaling a line by an element of \(F^\times\) makes
the ideal generated by its coordinates coincide with its chosen
representative \(\mathfrak a\). The subsequent ambiguity is exactly
\(\OO^\times\).

In positive orientation the principal Riemann form on \(E\) is
\(\operatorname{Im}(\bar zw)/\Delta\).
Pulling \(H\) back along \(\iota_v\) multiplies this form by \(Q(v)\).
A fundamental parallelogram for \(\mathfrak a^{-1}\) has area
\(\Delta/\NN(\mathfrak a)\), so \(\iota_v^*\Psi^*\mathcal L\) has
degree \(Q(v)/\NN(\mathfrak a)\). Since \(\iota_v\) is injective,
\(\ker(\Psi\iota_v)\) lies in \(E_{\mathfrak a}[n]\) and has order
\(k_{\mathfrak a}(\bar v)\). Dividing by the degree of this restricted
isogeny proves \eqref{eq:degree}.

To compare residue classes across ideal classes, choose
\(\lambda\in\mathfrak a\) generating \(\mathfrak a/n\mathfrak a\)
as a module over \(\OO/n\OO\). Such a generator exists by local
principality and the Chinese remainder theorem. Multiplication by
\(\lambda\) induces an isomorphism
\[
 \OO/n\OO \longrightarrow \mathfrak a/n\mathfrak a,\qquad
 E_{\mathfrak a}[n]\longrightarrow E[n].
\]
For the second isomorphism, observe that
\(\lambda\mathfrak a^{-1}\subset\OO\) has index coprime to \(n\).
For \(\bar v=\lambda\bar u\), the torsion equations
\(\iota_v(P)\in K\) and \(\iota_u(\lambda P)\in K\) coincide.
Thus the kernel orders and the admissible residues have the same
joint distribution in every ideal class. When \(n=1\), use the
single residue and kernel order one.

Fix an admissible residue \(\bar v\) in an ideal class
\(\mathfrak a\). Write \(k_0=k_{\mathfrak a}(\bar v)\) and
\(B=\{z\in\C^r:Q(z)\le1\}\). We count nonzero vectors in the residue
and in \(\sqrt{X\NN(\mathfrak a)k_0}\,B\). Primitivity is detected by
the M\"obius function \(\mu_F\) of integral ideals. For an ideal
\(\mathfrak b\) coprime to \(n\OO\), the intersection with
\((\mathfrak b\mathfrak a)^r\) is a coset of
\((n\mathfrak b\mathfrak a)^r\). The number of its points is
\begin{equation}\label{eq:latt}
 \frac{\vol(B)(Xk_0)^r}{n^{2r}\Delta^r\NN(\mathfrak b)^r}
       +O_{F,H,n}\bigl((X/\NN(\mathfrak b))^{r-1/2}+1\bigr).
\end{equation}
For the error, every ideal lattice \(\mathfrak b\mathfrak a\) is a
complex homothety of one of finitely many fixed representatives.
The rotation and scale act diagonally on \(\C^r\). Therefore a
fundamental cell has diameter bounded by a constant times
\(\sqrt{\NN(\mathfrak b\mathfrak a)}\) and volume
\(\bigl(\Delta\NN(\mathfrak b\mathfrak a)\bigr)^r\), before the fixed
scaling by \(n\). The cells meeting the ellipsoid boundary lie in a
boundary layer of that thickness. After normalization this gives
\eqref{eq:latt}, uniformly for all translates. Omitting the origin
changes the error by at most one.

For a nonzero vector counted by such an ideal, some coordinate
satisfies \(0\ne v_i\in\mathfrak b\mathfrak a\).
Then \(\NN(\mathfrak b\mathfrak a)\le\abs{v_i}^2\ll_{H,n}X\NN(\mathfrak a)\).
Hence the sum may stop at \(\NN(\mathfrak b)\ll_{H,n}X\).
There are \(O_F(Y)\) integral ideals of norm at most \(Y\).
For instance, in each ideal class they correspond, up to units, to
elements of a fixed fractional ideal of squared absolute value
bounded by a constant times \(Y\).
As \(r-1/2>1\), summing the errors in \eqref{eq:latt} gives
\(O_{F,H,n}(X^{r-1/2})\). Extending the main sum to infinity costs
\(O_{F,H,n}(X)\) and contributes
\[
 \sum_{(\mathfrak b,n)=1}\frac{\mu_F(\mathfrak b)}{\NN(\mathfrak b)^r}
 =
 \frac{1}{\zeta_F(r)\prod_{\mathfrak p\mid n\OO}(1-\NN(\mathfrak p)^{-r})}.
\]
The number of admissible residues in one class is
\(n^{2r}\prod_{\mathfrak p\mid n\OO}(1-\NN(\mathfrak p)^{-r})\).
Their sum therefore produces the indicated average. The unit action
on each nonzero vector is free, and multiplication by a unit leaves
both its norm and kernel order invariant. Divide by \(w_F\), sum
over the \(h_F\) classes and use
\(\vol(B)=\pi^r/(r!\det H)\). Finally, the primary decompositions of
\(K\) and \(E[n]\) factor the kernel order over rational primes.
Admissibility factors over the same primes, so the uniform average
factors.
\end{proof}

\begin{corollary}\label{cor:surface}
For a genus-two Jacobian with a CM-square isogeny presentation as
above, \eqref{eq:main} counts its geometric optimal quotient classes.
In the two-torsion graph case \(K=\operatorname{graph}(\gamma)\)
on \(E^2\), with \(H=2I\), the only finite calculation is
\begin{equation}\label{eq:graph}
 k(a,b)=\#\{t\in\OO/2\OO:\ bt=\gamma(at)\}.
\end{equation}
Here a choice of \(\OO/2\OO\)-generator identifies \(E[2]\) with
\(\OO/2\OO\), and \(\gamma\) denotes its additive action in those
coordinates.
\end{corollary}

For the geometric assertion, the polarized dual of an elliptic
inclusion is a quotient with connected kernel. Composing it with
an Abel--Jacobi morphism gives the optimal map. Conversely the
pullback of an optimal map is an elliptic inclusion, and
\(f_*f^*=[\deg f]\) identifies its degree with the polarization
degree. Changing the base point affects the induced curve map only
by a translation. Equation \eqref{eq:graph} then follows directly
from the graph condition. Changes of generator only relabel the
finite computation.

\section{Two CM Jacobians with common rational data}\label{sec:minus7}

Let \(F=\Q(\sqrt{-7})\) and
\[
  E:\quad Y^2=u^3-35u+98.
\]
Its \(j\)-invariant is \(-3375\), the \(j\)-invariant of the maximal
order of discriminant \(-7\) \cite{Cox}. The class number is one
by the imaginary quadratic Minkowski bound. Since \(2\) splits,
\(\OO/2\OO=\mathbb F_2\times\mathbb F_2\). Label the nonzero
elements \(e_1=(1,0)\), \(e_2=(0,1)\), \(s=e_1+e_2\).
The points \(e_1,e_2\) on \(E[2]\) have horizontal coordinates
\[
 a_1=(7+\sqrt{-7})/2,\qquad a_2=(7-\sqrt{-7})/2
\]
in either order, and \(s\) has coordinate \(a_3=-7\).
In fact complex conjugation acts on \(\OO\) by field conjugation and
interchanges the split idempotents, hence interchanges the two
eigenlines and fixes their sum.

Put \(\gamma_A=(e_1\ e_2)\) and \(\gamma_B=(e_2\ s)\).
Each graph is maximal isotropic for the product pairing on two-torsion.
The only automorphisms of \(E\) preserving the origin are \(\pm1\),
which act trivially on \(E[2]\). Neither of the two graphs is an
automorphism graph, so the corresponding principally polarized
quotients with pullback \(2I\) are genus-two Jacobians
\cite[Proposition 3]{HLP}.

\begin{proposition}\label{prop:minus7}
The two resulting genus-two curves \(C_A,C_B\) have isomorphic
polarized rational Tate modules at every prime over a common number
field, and equal Frobenius polynomials at its common good places.
Their geometric quotient counts satisfy
\[
 N_{C_A}(X)\sim \frac{3}{7L(2,\chi_{-7})}X^2,\qquad
 N_{C_B}(X)\sim \frac{4}{7L(2,\chi_{-7})}X^2 ,
\]
where \(\chi_{-7}\) is the quadratic Dirichlet character of \(F\).
\end{proposition}

\begin{proof}
The following three rows for each graph are indexed by \(a=e_1,e_2,s\).
A dash marks a nonadmissible pair. The row for \(a=0\) has one
admissible entry, \(b=s\), with kernel order one.
\begin{equation}\label{eq:split-table}
\begin{array}{c|cccc|cccc}
 &\multicolumn{4}{c|}{\gamma_A}&\multicolumn{4}{c}{\gamma_B}\\
 a\backslash b&0&e_1&e_2&s&0&e_1&e_2&s\\ \midrule
 e_1 & -&-&2&1&-&-&1&2\\
 e_2 & -&2&-&1&-&1&-&2\\
 s   & 1&1&1&2&1&2&2&2
\end{array}
\end{equation}
Each entry follows by testing the four \(t\) in \eqref{eq:graph}
with coordinatewise multiplication. Equivalently the histograms
are \(6,3\) and \(4,5\) for orders \(1,2\), respectively.
Thus \(\av k^2=2\) and \(8/3\). In \eqref{eq:constant} use
\(w_F=2\), \(\Delta^2=7/4\), \(\det H=4\) and
\(\zeta_F(2)=\zeta(2)L(2,\chi_{-7})\).

Over a common field on which the isogenies are defined, both
polarizations pull back to the same \(2I\). On rational Tate modules
the isogenies are isomorphisms, respect the pairings with that
common source polarization, and commute with Galois. This also
gives equality of the Frobenius polynomials at the stated places.
\end{proof}

The spectra already differ at degree three. A slope contributing
degree three must have total norm three and \(k=2\).
The norm \(a^2+ab+2b^2\) has two elements of norm one,
four of norm two, and none of norm three, as follows by completing
the square. The norm-one residues are \(s\), and the norm-two
residues are \(e_1,e_2\), two of each. Table~\eqref{eq:split-table}
gives order one for all such slopes for \(\gamma_A\) and order two
for all such slopes for \(\gamma_B\). Counting the two choices
of coordinate order and dividing by the two units gives
\(m_{C_A}(3)=0\) and \(m_{C_B}(3)=8\).

Here are equations for both curves and both maps in each case.
Let \(b_i=a_{\gamma(i)}\) with labels \(1,2,3\) as above. Put
\[
 q=\frac{b_2-b_1}{b_3-b_1},\qquad
 z=\frac{q(a_1-a_3)a_2-(a_1-a_2)a_3}{(a_1-a_2)-q(a_1-a_3)},
 \qquad r_i=a_i+z.
\]
Define \(\lambda=(b_2-b_1)/(1/r_2-1/r_1)\),
\(\mu=b_1-\lambda/r_1\), and choose \(\eta\) with
\(\eta^2=\mu^3-35\mu+98\). For \(\gamma_A,\gamma_B\) the
denominators are nonzero, as are \(r_i,\lambda,\eta\). The equation
and maps are
\begin{equation}\label{eq:models}
 C_\gamma:\quad y^2=\prod_i(x^2-r_i),\qquad
 f_1=(x^2-z,y),\qquad f_2=(\mu+\lambda/x^2,\eta y/x^3).
\end{equation}
Indeed, the defining equations give \(\mu+\lambda/r_i=b_i\).
Thus \(x^6\bigl((\mu+\lambda/x^2)^3-35(\mu+\lambda/x^2)+98\bigr)\)
has the three indicated roots in \(x^2\) and leading coefficient
\(\eta^2\). This verifies both map equations. The six branch
points are distinct, giving smooth genus two. Involution
\((x,y)\mapsto(-x,y)\) acts with opposite signs on the two elliptic
factors on Jacobians, so the pullbacks are orthogonal. The
pullbacks of two-torsion points at \(a_i\) and \(b_i\) coincide:
both give the sum of the two Weierstrass points over \(x^2=r_i\)
minus a hyperelliptic fibre. The graph therefore lies in the kernel
of \(f_1^*+f_2^*\), and both groups have order four. This verifies
the asserted gluing for the written models. For \(\gamma_A\) one
obtains \(z=5/3\) and the rational sextic
\(x^6-5x^4-\frac{80}{3}x^2+\frac{4096}{27}\).

\section{The full Gaussian multiplicity formula}\label{sec:gaussian}

Take \(F=\Q(i)\), \(E:y^2=x^3-x\), and \(H=2I\).
The images of the origin-preserving automorphisms in the permutation
group of the three nonzero two-torsion points form a subgroup
\(U\) of order two. If \(\gamma\in U\), the gluing comes from an
automorphism and the principally polarized quotient is a product
\cite{HLP}. The other four permutations form one double coset under
\(U\). Acting by an automorphism on either factor preserves \(2I\),
so all four non-product gluings have the same geometric spectrum.

Choose a generator of \(E[2]\) over \(R=\Z[i]/2\Z[i]\), and use the
permutation interchanging \(1\) and \(1+i\), with \(i\) fixed.
Let \(\chi_4(p)=1\) for primes \(p\equiv1\bmod4\) and
\(\chi_4(p)=-1\) for primes \(p\equiv3\bmod4\). For a positive odd
integer \(m\) set
\[
 B(m)=m\prod_{p\mid m}\left(1-\frac{\chi_4(p)}{p}\right),\qquad B(1)=1.
\]

\begin{theorem}\label{thm:gauss}
For the non-product Gaussian gluing the geometric multiplicities
are, for all odd \(m\),
\[
\begin{array}{c|c}
 d & m_C(d)\\ \midrule
 m,\ m\equiv3\bmod4 & 2B(m)\\
 m,\ m\equiv1\bmod4 & 0\\
 2m & 2(1+\mathbf 1_{m\equiv1\bmod4})B(m)\\
 4m & 2B(m)\\
 8\mid d & 0 .
\end{array}
\]
In particular there are four classes of degree two and eight of
degree three. If \(G=L(2,\chi_4)\) is Catalan's constant, then
\begin{equation}\label{eq:gauss-asympt}
 N_C(X)=\frac{15}{32G}X^2+O(X\log(2+X)).
\end{equation}
\end{theorem}

\begin{proof}
The ring \(R\) has units \(1,i\) and nonunits \(0,1+i\).
The admissible pairs have at least one unit. Substituting the four
possible \(t\) into \eqref{eq:graph} gives order two for equal
unit residues, or when a unit is paired with \(1+i\). The other
pairs have order one. In particular, the histogram is six of
each order among the twelve admissible pairs.

Let \(P(n)\) count primitive Gaussian pairs \((a,b)\) of total norm
\(n=|a|^2+|b|^2\), modulo the four units.
Jacobi's four-square formula gives \(8\sum_{d\mid n,\,4\nmid d}d\)
without the primitivity condition or the unit quotient
\cite{HardyWright}. If a Gaussian ideal of norm \(j\) divides both
coordinates, dividing by any of its generators scales the norm
by \(j\). Ideal M\"obius inversion now expresses \(P(n)/2\) as the
Dirichlet convolution of \(\sum_{d\mid n,\,4\nmid d}d\) with
\(b(j)=\sum_{\NN(\mathfrak a)=j}\mu_F(\mathfrak a)\).
The local generating series, with \(t\) marking the exponent of
the indicated odd prime \(p\), are
\[
 \frac{1-\chi_4(p)t}{1-pt}.
\]
At \(2\) the series is \(1+2t\): the initial four-square
divisor factor is \((1+2t)/(1-t)\) and the ideal M\"obius factor
is \(1-t\). Consequently
\[
 P(m)=2B(m),\qquad P(2m)=4B(m),\qquad P(4m)=0
 \quad(m\text{ odd}),
\]
and \(P(n)=0\) for every multiple of four.

For odd \(n\) exactly one coordinate is a unit modulo \(2\).
The other is \(1+i\) when \(n\equiv3\bmod4\), giving degree \(n\).
It is zero when \(n\equiv1\bmod4\), giving degree \(2n\).
For \(n\equiv2\bmod4\) both coordinates are units modulo \(2\).
Multiplying just one by \(i\) preserves primitiveness and norm and
induces a bijection between the equal and unequal residues, even
after dividing by common units. Half these classes have degree
\(n\) and half degree \(2n\). Combining the contributions gives
the table.

Finally \(B(m)=\sum_{j\mid m}\mu(j)\chi_4(j)\,m/j\).
For either odd residue \(a\bmod4\), sum first over odd \(j\) and
then over \(u=m/j\equiv aj^{-1}\bmod4\). The sum of positive
integers \(u\le Y\) in that residue is \(Y^2/8+O(Y+1)\). This
yields
\[
 \sum_{\substack{m\le X\\m\equiv a\bmod4}} B(m)
   =\frac{X^2}{8G}+O(X\log(2+X)).
\]
The four nonzero contributions from \(d=m,2m,4m\) have combined
coefficient \(15/(32G)\). This also agrees with
\eqref{eq:constant}, whose local average is \(5/2\) in this case.
\end{proof}

\section{Coprime spectra and models over the field of moduli}
\label{sec:twins}

Equality of counts after deleting the degrees divisible by one prime
can occur for inequivalent refined Humbert lattices. The case of level
three admits an exact explanation that works for every retained degree.

\begin{proposition}
Let \(G\) be an integral ternary lattice whose nondegenerate reduction
modulo \(3\) has square determinant. Its isotropic conic has four
projective points. Let \(L\) and \(L'\) be the inverse images of two
opposite secants, where opposite means that their endpoints partition
the four isotropic points. If \(3\nmid m\), the counts of ordinary and
primitive vectors of norm \(m^2\) agree in \(L\) and \(L'\).
\end{proposition}

\begin{proof}
The reduction is isometric over \(\mathbb F_3\) to
\(x^2+y^2+z^2\), by the classification of nondegenerate forms over
a finite field of odd characteristic. The conic is the four points
\([\,\pm1:\pm1:1\,]\). In this model, a vector of nonzero square norm
reduces to a coordinate axis. Each secant contains just one such
axis, and opposite secants contain the same one. Every vector of
norm \(m^2\), \(3\nmid m\), therefore lies in \(L\) exactly when it
lies in \(L'\). If a vector of norm \(m^2\) has content \(a\) in an
integral lattice basis, then \(a\mid m\). Divisor inversion gives the
primitive assertion.
\end{proof}

For \(D=-19\) two such lattices have Gram matrices, up to integral change
of basis,
\[
 Q_1=\begin{pmatrix}4&-2&0\\-2&12&-4\\0&-4&17\end{pmatrix},
 \qquad
 Q_2=\begin{pmatrix}4&0&2\\0&5&-2\\2&-2&36\end{pmatrix}.
\]
Their determinant is \(684\). The first difference is at \(m=3\),
where the primitive counts are \(0\) and \(4\). Their \(m=2\) counts
are \(2\) each. Both satisfy the realization criterion for primitive CM
refined Humbert forms \cite{KaniCM}. They do not represent one, so
the realizing principally polarized surfaces are Jacobians. Kani's
subcover correspondence identifies the primitive counts with geometric
optimal classes \cite{KaniI}. The opposite-secant proof is independent
of period recognition. The exact lattice and model data are recorded in the companion repository.

\subsection{An explicit descent for the bielliptic models}

Let
\[
 f(x)=x^6-s_1x^4+s_2x^2-1,\qquad
 u=s_1^3+s_2^3,\quad v=s_1s_2,\quad \Delta=u^2-4v^3.
\]
Suppose \(f\) is squarefree, \(v\Delta\ne0\), and \(u,v\in F\), where
\(F\) has characteristic zero. A model over \(F\) is
\begin{equation}\label{eq:fieldmodel}
 Y^2=aX^6+bX^5+cX^4+dX^3+\Delta cX^2+\Delta^2 bX+\Delta^3a,
\end{equation}
where
\begin{align*}
 a&=u-v^2,& b&=3\Delta+3u^2-2uv^2,\\
 c&=\Delta(15u+v^2),& d&=2\Delta(5\Delta+5u^2+2uv^2).
\end{align*}
Here model means a curve geometrically isomorphic to \(y^2=f(x)\).
The existence of a field-of-moduli model for the extra-involution
case belongs to Cardona and Quer \cite{CardonaQuer}. The formula
above gives a direct change of variables for the archived sextics.

Indeed, put \(s=s_1^3,\ t=s_2^3,\ \delta=s-t\) and write
\(P=X+\delta,\ Q=X-\delta\). Equation~\eqref{eq:fieldmodel}
is the binary sextic
\[
 \delta^{-1}
 \left(s^2P^6-sv^2P^4Q^2+tv^2P^2Q^4-t^2Q^6\right).
\]
It is invariant under simultaneous interchange \(s\leftrightarrow t\)
and \(P\leftrightarrow Q\) after division by \(\delta\). Expansion,
using \(\delta^2=\Delta\), gives the displayed coefficients. Over
an algebraic closure make the substitutions
\[
 x=\frac{s_1^2}{v}\frac{P}{Q},\qquad
 Y=\frac{t}{\sqrt{\delta}}Q^3y.
\]
Since \(st=v^3\), direct multiplication gives
\(t^2Q^6f(x)=s^2P^6-sv^2P^4Q^2+tv^2P^2Q^4-t^2Q^6\).
This proves the asserted isomorphism and smoothness.

\subsection{The recorded pair}

For the recorded \(D=-19\) models, both descended equations use
\[
 F=\Q[\theta],\qquad \theta^4-\theta^3-3\theta^2+2\theta+4=0 .
\]
For example, the two values of \(v\), in order \(Q_1,Q_2\), are
\[
\begin{aligned}
 &34441728\theta^3+4808832\theta^2-114871680\theta-100220535,\\
 &-1693056\theta^3-4394880\theta^2-3681408\theta-825975 .
\end{aligned}
\]
The companion data give the two \(u\) values and expanded sextics.
The replay checks both field embeddings and the coefficient identity
over the original degree-\(24\) field. Independently, on the original
sextic the Clebsch-normalized absolute ratio \(I_2^5/I_{10}\) is
\[
 \frac{16^5(v+15)^5}{64(v^2-4u+18v-27)^2}.
\]
Here \(I_2=16(v+15)\) is the degree-two transvectant invariant, and
the denominator is the discriminant of \(f\). The replay verifies
that the ratio agrees with the recorded absolute invariant after
the normalization factor is applied. Its two minimal polynomials
are irreducible quartics and differ. Thus the models' moduli fields
are quartic and their moduli points are in distinct Galois orbits.
The unordered pair of \(j\)-invariants of their two degree-two
elliptic quotients agrees exactly. The monic quadratic with those
two roots divides the conductor-three class polynomial
\(H_{-171}\) over \(F\).

\subsection{Certified identification with the secant lattices}

Write \(f(x)=g(x^2)\) and \(g^\vee(U)=-U^3g(1/U)\). The maps
\((x,y)\mapsto(x^2,y)\) and \((x,y)\mapsto(x^{-2},y/x^3)\) land on
\(Y^2=g(U)\) and \(Y^2=-g^\vee(U)\), respectively. For each root
\(\rho\) of \(g\), the pullbacks of \((\rho,0)\) and \((\rho^{-1},0)\)
as divisor classes coincide. Their difference as divisors is
\(\operatorname{div}(x)\), since they share the two Weierstrass
points above \(\rho\) and subtract the fibres over \(x=\infty\)
and \(x=0\), respectively. The Jacobian is therefore the quotient
of the product by the graph of this bijection on two-torsion.
Its polarization pulls back to twice the product polarization
\cite{HLP,KaniII}.

For each embedding of the equation field, exact divisibility by
\(H_{-171}\) restricts the elliptic \(j\)-invariants to the four
conductor-three CM classes. Certified complex enclosures separate
those four values and identify the two quotient lattices. To label
their nonzero two-torsion, let \(G_2,G_3\) be the Weierstrass
invariants of one normalized lattice and \(g_2,g_3\) those of the
centred cubic in Weierstrass form. The scaling \(\lambda\) satisfies
\(g_2=\lambda^{-4}G_2\) and \(g_3=\lambda^{-6}G_3\), so
\(\lambda^{-2}=g_3G_2/(G_3g_2)\). All these invariants are nonzero
for the lattices in question. Each scaled half-period value must
match a root of the centred cubic. Disjoint enclosures exclude five
of the six bijections, leaving the unique matching.

Adjoining lifts of the resulting graph gives the integral Jacobian
lattice. The replay verifies the descended polarization and computes
its refined Humbert form using rational arithmetic. It checks integral
change-of-basis matrices to \(Q_1\) and \(Q_2\), respectively, at all
\(24\) embeddings of the original equation field. The exact descent
then identifies the quartic models with the two secants. Thus their
geometric optimal quotient counts agree in every degree prime to
\(3\), and are \(0\) and \(4\), respectively, at degree \(3\).

\subsection{A second certified pair at discriminant \texorpdfstring{\(-43\)}{-43}}

The companion data give two further sextics over the quartic field
\(F_{-43}=\Q(\theta)\), where \(\theta^4-\theta^3+6\theta^2-\theta+7=0\).
They use the descent formula in \eqref{eq:fieldmodel}. For each, the
absolute invariants are checked in \(F_{-43}\), and the two elliptic
quotient \(j\)-invariants satisfy the conductor-three polynomial
\(H_{-387}\) by exact substitution. The same Arb and integral lattice
calculations identify the equations at all four embeddings with the
opposite-secant pair denoted \texttt{p3} in the data. The first difference
is at degree six, with zero and four geometric optimal classes. Both
curves have two degree-two classes. Equality in every degree prime to
three follows from the opposite-secant proposition.

For the six remaining discriminant \(-43\) lattices, the companion
records Mestre sextics over quadratic extensions of \(F_{-43}\) and
checks their absolute invariant triples exactly. The correspondence
between those six triples and the polarized lattices is based on
numerical recognition. The six models are outside the spectrum claim
for explicit equations above.

\section{Exact verification and scope}

The companion replay checks the two maps in \eqref{eq:models}
symbolically and evaluates the full finite kernel tables in both
quadratic rings. It then enumerates primitive vectors by Euclidean
division, with an integer norm cutoff. This is independent of the
four-square convolution. A nonprincipal ideal of norm two in
\(\Q(\sqrt{-5})\) tests the residue transfer used in the proof of
Theorem~\ref{thm:height}. For the Gaussian case the replay compares
every degree through \(150\) with both previously certified
N\'eron--Severi lattices. All three routes give \(11{,}574\)
geometric classes at \(94\) degrees.

The degree law assumes a specified finite kernel on the product
isogeny. From a curve equation, an independent route computes the
integral N\'eron--Severi lattice from verified elliptic maps or
correspondences, certified period enclosures and a verified bound
on the endomorphism rank. Its elliptic classes are the primitive
isotropic classes of positive polarization degree, as in
\cite{KaniI,KaniII}. On nine models the bounded homology search
checks \(1{,}204\) quotient directions individually against exact
lattices. The quaternionic model has a zero spectrum, which can
also be read from its N\'eron--Severi form by reduction modulo \(9\).
These calculations use the established endomorphism-ring framework
of \cite{CMSV,Lombardo}. Discovery of all the algebraic inputs of
that computation from an arbitrary curve equation is a separate
algorithmic problem.

For a quick reproducibility check, run the following from the
repository root with Python and SymPy installed.
\[
\begin{gathered}
\text{\ttfamily python -B research/elliptic-quotients/cm-spectra/verify.py}
\end{gathered}
\]
Run \texttt{python -B descend.py --verify} from the
\href{https://github.com/towre676-cloud/crest-root-calculus/tree/main/research/elliptic-quotients/twin-curves}{twin-curves directory}
to check the quartic models, then run
\texttt{python -B certify\_link.py --verify} with python-flint and
cypari2 to reproduce their Humbert forms. Run
\texttt{python -B verify\_d43.py} in the same directory to certify the
second pair. The
\href{https://github.com/towre676-cloud/crest-root-calculus/tree/main/docs/visuals}{visual companion}
shows the lattice-to-curve computation and the certificate in reverse.

The code, the exact lattice data, and the \TeX{} source are available
in the \href{https://github.com/towre676-cloud/crest-root-calculus}
{Crest--Root Calculus repository}.

\begin{thebibliography}{99}
\raggedright
\bibitem{CardonaQuer}
G.~Cardona and J.~Quer,
\emph{Field of moduli and field of definition for curves of genus 2},
in \emph{Computational Aspects of Algebraic Curves}, Lecture Notes Ser.
Comput. \textbf{13}, World Scientific, 2005, 71--83.
\href{https://arxiv.org/abs/math/0207015}{arXiv:math/0207015}.

\bibitem{CMSV}
E.~Costa, N.~Mascot, J.~Sijsling, and J.~Voight,
\emph{Rigorous computation of the endomorphism ring of a Jacobian},
Math. Comp. \textbf{88} (2019), 1303--1339.
\href{https://doi.org/10.1090/mcom/3373}{doi:10.1090/mcom/3373}.

\bibitem{Cox}
D.~A.~Cox, \emph{Primes of the form \(x^2+ny^2\)}, 2nd ed.,
John Wiley \& Sons, Hoboken, 2013.

\bibitem{Guerra19}
L.~Guerra, \emph{On elliptic curves of bounded degree in a polarized
Abelian variety}, Rend. Istit. Mat. Univ. Trieste \textbf{51}
(2019), 105--123.
\href{https://doi.org/10.13137/2464-8728/27068}{doi:10.13137/2464-8728/27068}.

\bibitem{HardyWright}
G.~H.~Hardy and E.~M.~Wright,
\emph{An introduction to the theory of numbers}, 6th ed.,
Oxford University Press, Oxford, 2008.

\bibitem{HLP}
E.~W.~Howe, F.~Lepr\'evost, and B.~Poonen,
\emph{Large torsion subgroups of split Jacobians of curves of genus
two or three}, Forum Math. \textbf{12} (2000), 315--364.
\href{https://arxiv.org/abs/math/9809210}{arXiv:math/9809210}.

\bibitem{KaniI}
E.~Kani, \emph{Elliptic subcovers of a curve of genus 2.
I. The isogeny defect}, Ann. Math. Qu\'e. \textbf{43} (2019), 281--303.
\href{https://doi.org/10.1007/s40316-018-0105-6}{doi:10.1007/s40316-018-0105-6}.

\bibitem{KaniII}
E.~Kani, \emph{Elliptic subcovers of a curve of genus 2.
II. The refined Humbert invariant},
J. Number Theory \textbf{193} (2018), 302--335.
\href{https://doi.org/10.1016/j.jnt.2018.05.011}{doi:10.1016/j.jnt.2018.05.011}.

\bibitem{KaniCM}
E.~Kani, \emph{The Refined Humbert Invariant for Abelian Product
Surfaces with Complex Multiplication},
\href{https://mast.queensu.ca/~kani/papers/refHum5.pdf}{manuscript}.

\bibitem{Lombardo}
D.~Lombardo, \emph{Computing the geometric endomorphism ring of a
genus 2 Jacobian}, Math. Comp. \textbf{88} (2019), 889--929.
\href{https://doi.org/10.1090/mcom/3358}{doi:10.1090/mcom/3358}.

\bibitem{Schanuel}
S.~H.~Schanuel, \emph{Heights in number fields},
Bull. Soc. Math. France \textbf{107} (1979), 433--449.
\href{https://doi.org/10.24033/bsmf.1905}{doi:10.24033/bsmf.1905}.
\end{thebibliography}
\end{document}
